% SPDX-License-Identifier: Apache-2.0
\section{A Function, Inlined}
\label{sec:x-fn}

The step of a unit is one loop body, and a long one reads as a wall.
A function under \code{\#[lower]} takes a piece of it out under a
name. It is combinational, a few \code{let}s and a value; it is plain
Rust for the simulation; and it is inlined at every call in a
lowered unit of the same file, its parameters bound to the arguments
and its \code{let}s to expressions of their own. The lowering reads
the file for it, since a proc macro sees one item at a time. What it
can inline is what the step could have said in place, and nothing
else: no register, no port, no \code{when!}. The counter here keeps
its two pieces of arithmetic so, the Gray code of the count and
whether the count is about to wrap, and the netlist is the one the
step would have made with them written in.

\begin{figure}[htbp]
\centering
\input{fn_timing.tex}
\caption{The counter: \code{code} is the Gray code of \code{n}, from
a function; \code{wrap} is the step that will wrap it.}
\label{fig:x-fn}
\end{figure}

\lstinputlisting[caption={\code{ex\_fn.rs}. Run with \code{bazel run //lib/examples:ex\_fn}.},label={lst:x-fn}]{lib/examples/ex_fn.rs}

\lstinputlisting[style=ascii,caption={What \code{ex\_fn} printed when the build ran it: the counts, then the lowering, the functions written in.},label={lst:x-fn-out}]{out_fn.txt}
